\section{Restricting the source-register operators}
\label{sec:peps_source_register_operators}

Let $\mathcal P$ be the set of parties, and let
$R=((p_i,q_i))_{i=0}^{n-1}$ be an ordered list of pairs of distinct parties.
For each $i$, let $U_i,V_i$ be complex inner product spaces. The ordered
register list is
\begin{align}
    \mathcal L_R
    = ((p_0,U_0),(q_0,V_0),\ldots,(p_{n-1},U_{n-1}),(q_{n-1},V_{n-1})).
    \label{eq:peps_source_register_order}
\end{align}
An empty register list has memory space $\C$. Register memories are tensored
in the order~(\ref{eq:peps_source_register_order}), with a final scalar factor.

\begin{definition}[The word of local source-register operators]
    \label{def:peps_source_register_operator_word}
    \lean{TNLean.Lattice.PartyCircuit.Word.localEndpoint}
    \lean{TNLean.Lattice.PartyCircuit.Word.sources_localEndpoint}
    \lean{TNLean.Lattice.PartyCircuit.SourceInventory.endpointWord}
    \lean{TNLean.Lattice.PartyCircuit.SourceInventory.sources_endpointWord}
    \leanok
    \uses{def:peps_party_layout,def:peps_source_slots}
    % Source: 04-compression.tex, lines 279--309 and 565--588.
    For the local matrices used in random source replacement and party
    contraction in \cite[proof of Theorem~5.2]{OpenAI2026PolynomialPEPS}, let
    $A_i:U_i\to U_i$ and $B_i:V_i\to V_i$ be continuous linear maps. The
    source-free word $Z_R=Z_R(A,B)$ applies $A_i$ at $(p_i,U_i)$ and $B_i$ at
    $(q_i,V_i)$ in increasing order of $i\in\{0,\ldots,n-1\}$, acting as the
    identity on all other registers; for $R=\varnothing$,
    $Z_{\varnothing}=\Id_{\C}$.
\end{definition}

\begin{theorem}[Restriction away from every source owner]
    \label{thm:peps_source_register_operator_exclusion}
    \lean{TNLean.Lattice.PartyCircuit.SourceInventory.restrict_endpointWord_eq_nil_and_eval_eq_id}
    \leanok
    \uses{def:peps_source_register_operator_word,thm:peps_party_restriction}
    Let $f:\mathcal P\to\{0,1\}$ satisfy $f(p_i)=f(q_i)=0$ for all
    $i\in\{0,\ldots,n-1\}$ (including $n=0$). Then $\mathcal L_R|_f=\varnothing$
    and $Z_R|_f=\Id_{\C}$.
\end{theorem}

\begin{proof}\leanok
    \uses{def:peps_source_register_operator_word,thm:peps_party_restriction,
        thm:peps_party_factorization}
    By induction on $R$ via Definition~\ref{def:peps_source_register_operator_word},
    every party occurring in $Z_R$ belongs to $\{p_i,q_i\}_{i=0}^{n-1}$, on
    which $f=0$. By Theorem~\ref{thm:peps_party_restriction} and
    Theorem~\ref{thm:peps_party_factorization}, $Z_R|_f$ is a source-free word
    with empty participating-party set and empty boundary lists
    $\mathcal L_R|_f=\varnothing$, hence acts as $\Id_{\C}$.
\end{proof}


\begin{theorem}[Action on the original endpoint factors]
    \label{thm:peps_source_register_operator_evaluation}
    \lean{TNLean.Lattice.PartyCircuit.SourceInventory.eval_endpointWord_slotVector_tmul}
    \leanok
    \uses{def:peps_source_register_operator_word,def:peps_joint_source_vector}
    For any $u_i\in U_i$ and $v_i\in V_i$, in the register
    order~(\ref{eq:peps_source_register_order}),
    \begin{align}
        Z_R\left(\bigotimes_{i=0}^{n-1}(u_i\otimes v_i)\right)
        =\bigotimes_{i=0}^{n-1}(A_i u_i\otimes B_i v_i),
        \label{eq:peps_source_register_operator_evaluation}
    \end{align}
    with empty tensor products equal to $1\in\C$ when $n=0$.
\end{theorem}

\begin{proof}
    \leanok
    \uses{def:peps_source_register_operator_word,def:peps_joint_source_vector}
    By induction on $R$ using Definition~\ref{def:peps_source_register_operator_word}
    and Definition~\ref{def:peps_joint_source_vector}, the operations at slot $0$
    map $u_0\otimes v_0\mapsto A_0u_0\otimes B_0v_0$ while acting as the
    identity on the tail registers, where the induction hypothesis yields~(\ref{eq:peps_source_register_operator_evaluation}).
\end{proof}

\begin{theorem}[Coordinates of the original joint source vector]
    \label{thm:peps_joint_source_vector_coordinates}
    \lean{TNLean.Lattice.PartyCircuit.SourceInventory.slotVector_eq_sum_slotBasis}
    \lean{TNLean.Lattice.PartyCircuit.SourceInventory.slotBasis_repr_slotVector}
    \leanok
    \uses{def:peps_joint_source_vector,def:peps_source_slot_basis}
    Let $b_i,c_i$ be finite orthonormal bases of $U_i,V_i$ and let $e_x$ be
    their ordered joint source basis (Definition~\ref{def:peps_source_slot_basis}).
    For any $\eta_i\in U_i\otimes V_i$,
    \begin{align}
        \bigotimes_i\eta_i
         & =\sum_x\left(\prod_i
        \langle b_{i,x_i^{(1)}}\otimes c_{i,x_i^{(2)}},\eta_i\rangle\right)e_x.
        \label{eq:peps_joint_source_coordinates}
    \end{align}
\end{theorem}

\begin{proof}\leanok
    \uses{def:peps_source_slot_basis,def:peps_joint_source_vector}
    Expand each $\eta_i$ in the orthonormal basis $b_i\otimes c_i$ and distribute
    $\bigotimes_i$ over the finite sums (Definition~\ref{def:peps_joint_source_vector}
    and Definition~\ref{def:peps_source_slot_basis}) to obtain~(\ref{eq:peps_joint_source_coordinates}).
\end{proof}

\begin{definition}[The operator with specified orthonormal coordinates]
    \label{def:peps_operator_from_basis_matrix}
    \lean{TNLean.Lattice.PartyCircuit.basisMatrixOperator}
    \lean{TNLean.Lattice.PartyCircuit.basisMatrixOperator_repr}
    \leanok
    For an orthonormal basis $b=(b_j)_{j\in I}$ of a finite-dimensional complex
    inner product space and a matrix $M=(M_{ij})_{i,j\in I}$, define $T_b(M)$ by
    $T_b(M)b_j=\sum_{i\in I} M_{ij}b_i$, so $[T_b(M)]_b=M$.
\end{definition}

\begin{theorem}[Coordinates of the source-endpoint operator]
    \label{thm:peps_endpoint_operator_coordinates}
    \lean{TNLean.Lattice.PartyCircuit.SourceInventory.slotBasis_repr_endpointWord}
    \lean{TNLean.Lattice.PartyCircuit.SourceInventory.endpointMatrixWord}
    \lean{TNLean.Lattice.PartyCircuit.SourceInventory.slotBasis_repr_endpointMatrixWord}
    \leanok
    \uses{thm:peps_source_register_operator_evaluation,
        thm:peps_joint_source_vector_coordinates,def:peps_operator_from_basis_matrix}
    In the joint source basis $(e_x)$, the operator $Z_R$ has matrix entries
    \begin{align}
        [Z_R]_{x,y} & =\prod_i
        \langle b_{i,x_i^{(1)}},A_i b_{i,y_i^{(1)}}\rangle
        \langle c_{i,x_i^{(2)}},B_i c_{i,y_i^{(2)}}\rangle.
        \label{eq:peps_endpoint_coordinates}
    \end{align}
    In particular, for $A_i=T_{b_i}(L_i)$ and $B_i=T_{c_i}(H_i)$
    (Definition~\ref{def:peps_operator_from_basis_matrix}),
    $[Z_R]_{x,y}=\prod_i (L_i)_{x_i^{(1)},y_i^{(1)}}(H_i)_{x_i^{(2)},y_i^{(2)}}$.
\end{theorem}

\begin{proof}\leanok
    \uses{thm:peps_source_register_operator_evaluation,
        thm:peps_joint_source_vector_coordinates,def:peps_operator_from_basis_matrix}
    Evaluate $Z_R e_y$ by Theorem~\ref{thm:peps_source_register_operator_evaluation}
    and extract the $e_x$-coordinate via~(\ref{eq:peps_joint_source_coordinates})
    and Definition~\ref{def:peps_operator_from_basis_matrix}.
\end{proof}

\begin{theorem}[Finite averaging of endpoint operators]
    \label{thm:peps_endpoint_matrix_average}
    \lean{TNLean.Lattice.PartyCircuit.SourceInventory.endpointMatrixAverage_eq_product}
    \leanok
    \uses{thm:peps_endpoint_operator_coordinates}
    For endpoint matrices $L_{i,s},H_{i,s}$ ($0\le i<n$, $0\le s<k$), entrywise
    in the joint source basis,
    \begin{align}
        \bigotimes_{i=0}^{n-1}
        \left(k^{-1}\sum_{s=0}^{k-1} L_{i,s}\otimes H_{i,s}\right)
         & =k^{-n}\sum_{j:\{0,\ldots,n-1\}\to\{0,\ldots,k-1\}} Z_j,
        \label{eq:peps_endpoint_average}
    \end{align}
    where $Z_j=Z_R(T_b(L_{\cdot,j(\cdot)}),T_c(H_{\cdot,j(\cdot)}))$ (with $0^{-1}=0$ for $k=0$).
\end{theorem}

\begin{proof}\leanok
    \uses{thm:peps_endpoint_operator_coordinates}
    Substitute~(\ref{eq:peps_endpoint_coordinates}) from
    Theorem~\ref{thm:peps_endpoint_operator_coordinates} and distribute
    $\prod_{i=0}^{n-1}k^{-1}\sum_{s=0}^{k-1}$ over assignments
    $j:\{0,\ldots,n-1\}\to\{0,\ldots,k-1\}$.
\end{proof}

\begin{theorem}[Matrices in orthonormal coordinates]
    \label{thm:peps_matrix_orthonormal_coordinates}
    \lean{TNLean.Lattice.PartyCircuit.matrix_repr_eval_eq_toMatrix}
    \leanok
    For finite orthonormal bases $b=(b_j)$ and $c=(c_i)$ and a continuous linear
    map $A$, the matrix of $A$ has entries $M_{ij}=\langle c_i,A b_j\rangle$.
\end{theorem}

\begin{proof}\leanok
    Take the $c_i$-coordinate of $Ab_j$ by orthonormal expansion.
\end{proof}

\begin{theorem}[Dependence on the selected endpoint maps]
    \label{thm:peps_source_register_operator_local_congruence}
    \lean{TNLean.Lattice.PartyCircuit.SourceInventory.restrict_endpointWord_eval_eq_of_local_eq}
    \leanok
    \uses{def:peps_source_register_operator_word,thm:peps_party_restriction}
    % Source: 04-compression.tex, lines 279--309 and 565--588.
    For a party selector $f:\mathcal P\to\{0,1\}$ and endpoint maps $A_i,B_i$
    and $A'_i,B'_i$ in \cite[proof of Theorem~5.2]{OpenAI2026PolynomialPEPS}, if
    $A_i=A'_i$ whenever $f(p_i)=1$ and $B_i=B'_i$ whenever $f(q_i)=1$, then
    \begin{align}
        Z_R(A,B)|_f & =Z_R(A',B')|_f.
        \label{eq:peps_source_register_operator_local_congruence}
    \end{align}
\end{theorem}

\begin{proof}\leanok
    \uses{def:peps_source_register_operator_word,thm:peps_party_restriction}
    By induction on $R$ via Definition~\ref{def:peps_source_register_operator_word}
    and Theorem~\ref{thm:peps_party_restriction}, each unselected endpoint map is
    omitted under $(\cdot)|_f$ while each selected endpoint map satisfies
    $A_i=A'_i$ or $B_i=B'_i$ by hypothesis.
\end{proof}


\begin{theorem}[Restriction and register-list identification]
    \label{thm:peps_party_restriction_canonical_layouts}
    \lean{TNLean.Lattice.PartyCircuit.Word.restrict_castLayouts}
    \leanok
    \uses{thm:peps_party_restriction,thm:peps_basis_memory_equalities}
    % Source: 04-compression.tex, lines 233--251, 279--309 and 565--588.
    For a source-free chronological word $W:a\to b$, register-list equalities
    $a=a'$ and $b=b'$, and a party selector $f$ in
    \cite[proof of Theorem~5.2]{OpenAI2026PolynomialPEPS}, restricting $W$ after
    identifying $a=a'$ and $b=b'$ equals restricting $W|_f$ first and then
    identifying $a|_f=a'|_f$ and $b|_f=b'|_f$.
\end{theorem}

\begin{proof}\leanok
    \uses{thm:peps_party_restriction,thm:peps_basis_memory_equalities}
    Substituting the list equalities $a=a'$ and $b=b'$ via
    Theorem~\ref{thm:peps_basis_memory_equalities} commutes with the filter
    $(\cdot)|_f$ of Theorem~\ref{thm:peps_party_restriction}.
\end{proof}

\begin{theorem}[Dependence on incident source samples]
    \label{thm:peps_actual_sampled_endpoint_locality}
    \lean{TNLean.Lattice.PartyCircuit.SourceCircuit.restrict_sampledEndpointInputWord_eval_eq_of_incident_samples}
    \leanok
    \uses{def:peps_actual_sampled_endpoint_word,
        def:peps_original_sampled_endpoint_factors,
        thm:peps_source_register_operator_local_congruence,
        thm:peps_party_restriction_canonical_layouts}
    % Source: 04-compression.tex, lines 279--309 and 565--588.
    For ket/bra choices $\xi,\zeta$, a grouped-owner selector $f$, and two sample
    assignments $(j,\omega)$ and $(j',\omega')$ (Definition~\ref{def:peps_source_gaussian_law})
    in \cite[proof of Theorem~5.2]{OpenAI2026PolynomialPEPS}, let
    $U_{\xi,\zeta,j,\omega}$ be the sampled source-input word of
    Definition~\ref{def:peps_actual_sampled_endpoint_word}. If at every source
    $e$ with at least one endpoint selected by $f$ and branch pair
    $(\alpha,\beta)$ under $(\xi,\zeta)$,
    \begin{align}
        j(e)=j'(e)
        \qquad\text{and}\qquad
        \omega_{e,\alpha,\beta,s,p,q}=\omega'_{e,\alpha,\beta,s,p,q}
        \quad\text{for all }s,p,q,
        \label{eq:peps_actual_sampled_endpoint_locality_agreement}
    \end{align}
    then
    \begin{align}
        U_{\xi,\zeta,j,\omega}|_f
         & =U_{\xi,\zeta,j',\omega'}|_f.
        \label{eq:peps_actual_sampled_endpoint_locality}
    \end{align}
\end{theorem}

\begin{proof}\leanok
    \uses{def:peps_actual_sampled_endpoint_word,
        def:peps_original_sampled_endpoint_factors,
        thm:peps_source_register_operator_local_congruence,
        thm:peps_party_restriction_canonical_layouts}
    By Definition~\ref{def:peps_original_sampled_endpoint_factors},
    Definition~\ref{def:peps_actual_sampled_endpoint_word}, and~(\ref{eq:peps_actual_sampled_endpoint_locality_agreement}),
    the endpoint matrices agree at every endpoint selected by $f$.
    Applying Theorem~\ref{thm:peps_source_register_operator_local_congruence} and
    Theorem~\ref{thm:peps_party_restriction_canonical_layouts} yields~(\ref{eq:peps_actual_sampled_endpoint_locality}).
\end{proof}

\begin{theorem}[The physical mixed trace in the original coordinates]
    \label{thm:peps_selected_physical_private_trace}
    \lean{TNLean.Lattice.PartyCircuit.selectedPhysicalPrivateBasis_partialTrace_mixed}
    \leanok
    \uses{thm:peps_selected_physical_private_basis_grouped,
        thm:peps_selected_physical_basis_grouped,def:peps_selected_physical_configurations}
    % Source: 04-compression.tex, lines 137--151, 409--434 and 565--588.
    Let $H$ be the original physical memory and $G$ the retained private memory
    in \cite[proof of Theorem~5.2]{OpenAI2026PolynomialPEPS}, with grouped
    physical/private basis $B$, coordinate-separation map $\chi$ (removing the
    exterior singleton coordinate), and product readout $J$ with respect to the
    original physical basis and any finite orthonormal reference basis of $G$.
    For $v,t\in H\otimes G$ with grouped coordinates $b_v,b_t$,
    \begin{align}
        \operatorname{tr}_{\mathrm{private}}
        \left((b_vb_t^\dagger)^\chi\right)
         & =\operatorname{tr}_{G}\left((Jv)(Jt)^\dagger\right),
        \label{eq:peps_original_physical_mixed_trace}
    \end{align}
    indexed by $\prod_{p\in\mathcal P}\{0,\ldots,d_p-1\}$
    (Definition~\ref{def:peps_selected_physical_configurations}).
\end{theorem}

\begin{proof}\leanok
    \uses{thm:peps_selected_physical_private_basis_grouped,
        thm:peps_selected_physical_basis_grouped}
    By Theorem~\ref{thm:peps_selected_physical_private_basis_grouped} and
    Theorem~\ref{thm:peps_selected_physical_basis_grouped}, $\chi\circ B$
    coincides with the product of the original physical basis and the grouped
    private basis on $G$. Any unitary change of basis $I\otimes V$ on $G$ cancels
    under $\operatorname{tr}_G$, giving~(\ref{eq:peps_original_physical_mixed_trace}).
\end{proof}



\begin{theorem}[The original physical density in grouped coordinates]
    \label{thm:peps_distributed_grouped_physical_trace}
    \lean{TNLean.Lattice.PartyCircuit.DistributedConstruction.physicalDensity_eq_grouped_partialTrace}
    \leanok
    \uses{def:peps_distributed_construction_density,
        thm:peps_selected_physical_private_trace}
    For a distributed construction with circuit output $v$ and grouped
    coordinates $b_v$, in the notation of
    Theorem~\ref{thm:peps_selected_physical_private_trace} and
    Definition~\ref{def:peps_distributed_construction_density},
    $\rho=\operatorname{tr}_{\mathrm{private}}((b_vb_v^\dagger)^\chi)$.
    % Source: 04-compression.tex, lines 137--153 and 574--579.
\end{theorem}

\begin{proof}\leanok
    \uses{def:peps_distributed_construction_density,
        thm:peps_selected_physical_private_trace}
    Set $t=v$ in~(\ref{eq:peps_original_physical_mixed_trace}) of
    Theorem~\ref{thm:peps_selected_physical_private_trace} using the reference
    private basis of Definition~\ref{def:peps_distributed_construction_density}.
\end{proof}

\begin{theorem}[Restriction of an exchange of register blocks]
    \label{thm:peps_party_restriction_exchange_blocks}
    \lean{TNLean.Lattice.PartyCircuit.Word.restrict_exchangeBlocks_eval_heq}
    \leanok
    \uses{thm:peps_party_restriction,thm:peps_basis_memory_equalities}
    % Source: 04-compression.tex, lines 233--267, 409--434 and 565--588.
    For ordered register lists $a,b,t$, an owner selector $f$, and the block-exchange
    word $S_{a,b;t}:abt\to bat$ in
    \cite[proof of Theorem~5.2]{OpenAI2026PolynomialPEPS}, under the retained-list
    identifications,
    \begin{align}
        (S_{a,b;t})|_f
         & =S_{a|_f,b|_f;t|_f}.
        \label{eq:peps_restricted_block_exchange}
    \end{align}
\end{theorem}

\begin{proof}\leanok
    \uses{thm:peps_party_restriction,thm:peps_basis_memory_equalities}
    By Theorem~\ref{thm:peps_party_restriction} and
    Theorem~\ref{thm:peps_basis_memory_equalities}, moving a single register
    through $a$ restricts to moving it through $a|_f$ when retained and to $\Id$
    when erased; induction on $b$ then yields~(\ref{eq:peps_restricted_block_exchange}).
\end{proof}

\begin{theorem}[Restriction beside spectator registers]
    \label{thm:peps_append_spectators_restriction}
    \lean{TNLean.Lattice.PartyCircuit.Word.restrict_appendTail_eval_heq}
    \leanok
    \uses{def:peps_append_spectators,thm:peps_party_restriction,
        thm:peps_party_restriction_exchange_blocks,
        thm:peps_party_restriction_canonical_layouts}
    % Source: 04-compression.tex, lines 233--267, 409--434 and 565--588.
    For a source-free chronological word $W:a\to b$, a spectator list $t$, and
    an owner selector $f$ in \cite[proof of Theorem~5.2]{OpenAI2026PolynomialPEPS},
    the spectator extension $W_t$ (Definition~\ref{def:peps_append_spectators})
    satisfies $(W_t)|_f=(W|_f)_{t|_f}$ under the retained-list identifications.
\end{theorem}

\begin{proof}\leanok
    \uses{def:peps_append_spectators,thm:peps_party_restriction,
        thm:peps_party_restriction_exchange_blocks,
        thm:peps_party_restriction_canonical_layouts}
    By induction on $t$, restriction commutes with framing
    (Theorem~\ref{thm:peps_party_restriction}); combining this with~(\ref{eq:peps_restricted_block_exchange})
    and Theorem~\ref{thm:peps_party_restriction_canonical_layouts} on the three
    factors of $W_t$ in Definition~\ref{def:peps_append_spectators} gives
    $(W_t)|_f=(W|_f)_{t|_f}$.
\end{proof}

\begin{theorem}[Spectators at an excluded owner]
    \label{thm:peps_spectator_owner_exclusion}
    \lean{TNLean.Lattice.PartyCircuit.Word.restrict_mapOwner_appendTail_eval_heq_id_of_disjoint}
    \leanok
    \uses{thm:peps_owner_disjoint_restriction,
        thm:peps_append_spectators_restriction,thm:peps_append_spectators}
    % Source: 04-compression.tex, lines 233--267, 351--381, 409--434 and 565--588.
    Let $U:a\to b$ be a source-free chronological word with owners in $C$, let
    $\gamma:C\to\mathcal P$ satisfy $\gamma(c)\ne p$ for all $c\in C$, and let
    $t$ be a spectator list with owners in $\mathcal P$
    (\cite[proof of Theorem~5.2]{OpenAI2026PolynomialPEPS}). Then
    $((U^\gamma)_t)|_p=\Id_{\mathcal H_{t|_p}}$.
\end{theorem}

\begin{proof}\leanok
    \uses{thm:peps_owner_disjoint_restriction,
        thm:peps_append_spectators_restriction,thm:peps_append_spectators}
    By Theorem~\ref{thm:peps_owner_disjoint_restriction}, $(U^\gamma)|_p=\Id_{\C}$
    on empty boundary lists. Applying Theorem~\ref{thm:peps_append_spectators_restriction}
    and Theorem~\ref{thm:peps_append_spectators} gives
    $((U^\gamma)_t)|_p=((U^\gamma)|_p)_{t|_p}=\Id_{\C}\otimes\Id_{\mathcal H_{t|_p}}=\Id_{\mathcal H_{t|_p}}$.
\end{proof}

\begin{theorem}[Operator identities under composition and framing]
    \label{thm:peps_operator_composition_framing}
    \lean{TNLean.Lattice.PartyCircuit.Word.eval_comp_heq,
        TNLean.Lattice.PartyCircuit.Word.eval_frame_heq,
        TNLean.Lattice.PartyCircuit.Word.eval_frameList_heq}
    \leanok
    \uses{thm:peps_basis_memory_equalities,def:peps_party_registers}
    % Source: 04-compression.tex, lines 233--251 and 409--434.
    For register-list equalities $a=a'$, $b=b'$, $c=c'$ and $t=t'$ in
    \cite[proof of Theorem~5.2]{OpenAI2026PolynomialPEPS}, if chronological words
    $U:a\to b$, $V:b\to c$, $U':a'\to b'$, $V':b'\to c'$ satisfy $U=U'$ and
    $V=V'$ under the memory identifications, then $V\circ U=V'\circ U'$,
    $I_{\mathcal H_r}\otimes U=I_{\mathcal H_r}\otimes U'$ for any register $r$,
    and $F_t(U)=F_{t'}(U')$.
\end{theorem}

\begin{proof}\leanok
    \uses{thm:peps_basis_memory_equalities}
    By Theorem~\ref{thm:peps_basis_memory_equalities}, identifying the equal
    register lists reduces $V\circ U=V'\circ U'$ and
    $I_{\mathcal H_r}\otimes U=I_{\mathcal H_r}\otimes U'$ to congruence of
    composition and tensoring, and induction on $t=t'$ gives $F_t(U)=F_{t'}(U')$.
\end{proof}

\begin{theorem}[Restriction of the local part without sources]
    \label{thm:peps_source_free_local_part_restriction}
    \lean{TNLean.Lattice.PartyCircuit.Word.restrict_localPart_eval_heq_of_sources_nil}
    \leanok
    \uses{thm:peps_source_preparation,thm:peps_party_restriction,
        thm:peps_basis_memory_equalities}
    % Source: 04-compression.tex, lines 233--267, 351--381 and 409--434.
    For a chronological word $W:a\to b$ with empty source inventory
    $s_W=\varnothing$ and an owner selector $f$ in
    \cite[proof of Theorem~5.2]{OpenAI2026PolynomialPEPS}, the local part
    $L(W)$ (Theorem~\ref{thm:peps_source_preparation}) satisfies
    $(L(W))|_f=W|_f$ under the identification $\varnothing a=a$.
\end{theorem}

\begin{proof}\leanok
    \uses{thm:peps_source_preparation,thm:peps_party_restriction,
        thm:peps_excluded_spectator_restriction,
        thm:peps_operator_composition_framing,thm:peps_register_exchange}
    By induction on $W$ using Theorem~\ref{thm:peps_source_preparation} and
    Theorem~\ref{thm:peps_party_restriction}, identities, local operations, and
    exchanges satisfy $(L(W))|_f=W|_f$ directly, and source preparations are
    excluded by $s_W=\varnothing$. For compositions and framed words, each
    component has $s=\varnothing$, so moving through or framing by $s$ acts as
    $\Id$ (Theorem~\ref{thm:peps_excluded_spectator_restriction} and
    Theorem~\ref{thm:peps_register_exchange}), and
    Theorem~\ref{thm:peps_operator_composition_framing} closes the induction.
\end{proof}

\begin{theorem}[Moving a register through an erased prefix]
    \label{thm:peps_erased_prefix_register_move}
    \lean{TNLean.Lattice.PartyCircuit.Word.restrict_moveHead_eval_heq_id_of_prefix_eq_nil}
    \leanok
    \uses{thm:peps_party_restriction,thm:peps_register_exchange,
        thm:peps_basis_memory_equalities}
    % Source: 04-compression.tex, lines 233--267, 351--381, 409--434 and 565--588.
    Let $r$ be a register, let $s,t$ be ordered register lists, and let $f$ be
    an owner selector with $s|_f=\varnothing$
    (\cite[proof of Theorem~5.2]{OpenAI2026PolynomialPEPS}). Under the
    identification $(srt)|_f=(rst)|_f=(rt)|_f$, the register move
    $M_r(s,t):srt\to rst$ satisfies $(M_r(s,t))|_f=\Id_{\mathcal H_{(rt)|_f}}$.
\end{theorem}

\begin{proof}\leanok
    \uses{thm:peps_party_restriction,thm:peps_register_exchange,
        thm:peps_basis_memory_equalities}
    By Theorem~\ref{thm:peps_party_restriction} and
    Theorem~\ref{thm:peps_register_exchange}, since $s|_f=\varnothing$,
    $(M_r(s,t))|_f$ moves $r$ through $\varnothing$ if $r$ is retained (acting
    as $\Id_{\mathcal H_{r(t|_f)}}$) and acts as $\Id_{\mathcal H_{t|_f}}$ if
    $r$ is erased; Theorem~\ref{thm:peps_basis_memory_equalities} identifies
    both with $\Id_{\mathcal H_{(rt)|_f}}$.
\end{proof}

\begin{theorem}[Framing a local part with erased sources]
    \label{thm:peps_local_part_framing_restriction}
    \lean{TNLean.Lattice.PartyCircuit.Word.restrict_localPart_frameList_eval_heq}
    \leanok
    \uses{thm:peps_source_preparation,thm:peps_party_restriction,
        thm:peps_basis_memory_equalities}
    % Source: 04-compression.tex, lines 233--267, 351--381, 409--434 and 565--588.
    Let $W:a\to b$ be a chronological word with source-register list $s$, let
    $f$ satisfy $s|_f=\varnothing$, and let $t$ be a spectator list
    (\cite[proof of Theorem~5.2]{OpenAI2026PolynomialPEPS}). For the framed
    word $F_t(W)$ and local part $L(W):sa\to b$
    (Theorem~\ref{thm:peps_source_preparation}), under the retained-list
    identifications,
    \begin{align}
        (L(F_t(W)))|_f
         & =F_{t|_f}((L(W))|_f).
        \label{eq:peps_local_part_framing}
    \end{align}
\end{theorem}

\begin{proof}\leanok
    \uses{thm:peps_source_preparation,thm:peps_party_restriction,
        thm:peps_erased_prefix_register_move,
        thm:peps_operator_composition_framing,thm:peps_basis_memory_equalities}
    By induction on $t$, $t=\varnothing$ is immediate. For $t=rt'$,
    Theorem~\ref{thm:peps_source_preparation} expresses $L(F_{rt'}(W))$ as
    $F_r(L(F_{t'}(W)))\circ M_r(s,t'a)$. Since $s|_f=\varnothing$,
    $(M_r(s,t'a))|_f=\Id$ by Theorem~\ref{thm:peps_erased_prefix_register_move},
    and applying the induction hypothesis with
    Theorem~\ref{thm:peps_party_restriction} and
    Theorem~\ref{thm:peps_operator_composition_framing} yields~(\ref{eq:peps_local_part_framing}).
\end{proof}

\begin{theorem}[Identities and input identifications under restriction]
    \label{thm:peps_identity_input_identification}
    \lean{TNLean.Lattice.PartyCircuit.Word.eval_id_heq,
        TNLean.Lattice.PartyCircuit.Word.restrict_castInput_eval_heq}
    \leanok
    \uses{thm:peps_basis_memory_equalities,thm:peps_party_restriction}
    % Source: 04-compression.tex, lines 233--251 and 409--434.
    For an equality $a=a'$ of ordered register lists in
    \cite[proof of Theorem~5.2]{OpenAI2026PolynomialPEPS},
    $\Id_{\mathcal H_a}=\Id_{\mathcal H_{a'}}$ under the memory identification,
    and for any source-free word $W:a\to b$ and owner selector $f$, replacing
    $a$ by $a'$ before restricting equals $W|_f$ under $a|_f=a'|_f$.
\end{theorem}

\begin{proof}\leanok
    \uses{thm:peps_basis_memory_equalities,thm:peps_party_restriction}
    Substituting $a=a'$ via Theorem~\ref{thm:peps_basis_memory_equalities} and
    Theorem~\ref{thm:peps_party_restriction} reduces both claims to reflexivity.
\end{proof}

\begin{theorem}[Spectators outside the owners of a sourced composition]
    \label{thm:peps_sourced_spectator_owner_exclusion}
    \lean{TNLean.Lattice.PartyCircuit.Word.restrict_localPart_mapOwner_appendTail_eval_heq_id_of_disjoint}
    \leanok
    \uses{thm:peps_source_word_parties,def:peps_append_spectators,
        thm:peps_party_restriction,thm:peps_basis_memory_equalities}
    % Source: 04-compression.tex, lines 233--267, 351--381, 409--434 and 565--588.
    Let $U:a\to b$ be a chronological word with owners in $C$, let
    $\gamma:C\to\mathcal P$ satisfy $\gamma(c)\ne p$ for all $c\in C$, and let
    $t$ be a spectator list (\cite[proof of Theorem~5.2]{OpenAI2026PolynomialPEPS}).
    Under the retained-list identifications,
    \begin{align}
        \bigl(L((U^\gamma)_t)\bigr)|_p
         & =\Id_{\mathcal H_{t|_p}}.
        \label{eq:peps_sourced_spectator_exclusion}
    \end{align}
\end{theorem}

\begin{proof}\leanok
    \uses{thm:peps_source_word_parties,def:peps_append_spectators,
        thm:peps_source_free_local_part_restriction,
        thm:peps_party_restriction_exchange_blocks,
        thm:peps_party_restriction_canonical_layouts,
        thm:peps_local_part_framing_restriction,
        thm:peps_excluded_spectator_restriction,
        thm:peps_operator_composition_framing,
        thm:peps_identity_input_identification}
    Since $p\notin\gamma(C)$, $(a^\gamma)|_p=(b^\gamma)|_p=(s_{U^\gamma})|_p=\varnothing$,
    and $(L(U^\gamma))|_p=\Id_{\C}$ by Theorem~\ref{thm:peps_source_word_parties}.
    Expanding $(U^\gamma)_t$ by Definition~\ref{def:peps_append_spectators}, its
    initial and final block exchanges restrict to $\Id_{\mathcal H_{t|_p}}$ by
    Theorem~\ref{thm:peps_source_free_local_part_restriction} and~(\ref{eq:peps_restricted_block_exchange}),
    while the middle framed local part restricts to
    $F_{t|_p}(\Id_{\C})=\Id_{\mathcal H_{t|_p}}$ by~(\ref{eq:peps_local_part_framing})
    and Theorem~\ref{thm:peps_excluded_spectator_restriction}. Composing via
    Theorem~\ref{thm:peps_operator_composition_framing} and
    Theorem~\ref{thm:peps_identity_input_identification} gives~(\ref{eq:peps_sourced_spectator_exclusion}).
\end{proof}

\begin{theorem}[Spectators at a party absent from the core]
    \label{thm:peps_absent_party_spectator_identity}
    \lean{TNLean.Lattice.PartyCircuit.Word.restrict_localPart_appendTail_eval_heq_id_of_not_mem_parties}
    \leanok
    \uses{thm:peps_source_word_parties,def:peps_append_spectators,
        thm:peps_party_restriction,def:peps_single_party_selector}
    % Source: 04-compression.tex, lines 233--267, 351--381, 409--434 and 565--588.
    Let $U:a\to b$ be a chronological word, let $p\notin\mathcal P(U)$, and let
    $t$ be a spectator list (\cite[proof of Theorem~5.2]{OpenAI2026PolynomialPEPS}).
    Under the retained-list identifications,
    $(L(U_t))|_p=\Id_{\mathcal H_{t|_p}}$.
\end{theorem}

\begin{proof}\leanok
    \uses{thm:peps_source_word_parties,thm:peps_party_restriction,
        thm:peps_party_factorization,def:peps_append_spectators,
        thm:peps_source_free_local_part_restriction,
        thm:peps_party_restriction_exchange_blocks,
        thm:peps_party_restriction_canonical_layouts,
        thm:peps_local_part_framing_restriction,
        thm:peps_excluded_spectator_restriction,
        thm:peps_operator_composition_framing,
        thm:peps_identity_input_identification}
    Since $p\notin\mathcal P(U)$, $a|_p=b|_p=s_U|_p=\varnothing$, and
    $\mathcal P(L(U))=\mathcal P(U)$ by Theorem~\ref{thm:peps_source_word_parties},
    so $(L(U))|_p=\Id_{\C}$ by Theorem~\ref{thm:peps_party_restriction} and
    Theorem~\ref{thm:peps_party_factorization}. As in the proof of
    Theorem~\ref{thm:peps_sourced_spectator_owner_exclusion}, the three factors
    of $U_t$ (Definition~\ref{def:peps_append_spectators}) each restrict at $p$
    to $\Id_{\mathcal H_{t|_p}}$, and composing gives
    $(L(U_t))|_p=\Id_{\mathcal H_{t|_p}}$.
\end{proof}

\begin{theorem}[A nonincident party in the all-party expansion]
    \label{thm:peps_all_gate_nonincident_identity}
    \lean{TNLean.Lattice.PartyCircuit.SourceCircuit.restrict_partialWord_all_gate_localPart_eval_heq_id_of_nonincident}
    \leanok
    \uses{def:peps_prepared_source_gate,def:peps_partial_source_choices,
        def:peps_affected_owner,def:peps_single_party_selector}
    % Source: 04-compression.tex, lines 351--381, 409--434 and 565--588.
    Let $G$ be a prepared gate (Definition~\ref{def:peps_prepared_source_gate})
    with participant set $C$, branch set $\Xi$, and injection
    $\gamma:C\hookrightarrow P$, let $t$ be a spectator list, and fix $p\in P$
    with $p\notin\gamma(C)$ (\cite[proof of Theorem~5.2]{OpenAI2026PolynomialPEPS}).
    Let $g:P\to Q=P\sqcup\{*\}$ be the all-party owner map
    (Definition~\ref{def:peps_affected_owner}). For any partial choice $\xi$
    (Definition~\ref{def:peps_partial_source_choices}) with partial composition
    $V_\xi$, under the retained-list identifications,
    \begin{align}
        (L(V_\xi))|_{g(p)}
         & =\Id_{\mathcal H_{(t^g)|_{g(p)}}}.
        \label{eq:peps_all_gate_nonincident}
    \end{align}
\end{theorem}

\begin{proof}\leanok
    \uses{def:peps_partial_source_choices,def:peps_affected_owner,
        thm:peps_source_word_parties,thm:peps_owner_spectator_extension,
        thm:peps_absent_party_spectator_identity,
        thm:peps_source_free_local_part_restriction,
        def:peps_grouped_block_map,thm:peps_party_restriction,
        thm:peps_party_restriction_canonical_layouts}
    If $C\ne\varnothing$, commuting $g$ past the spectator extension by $t$ via
    Theorem~\ref{thm:peps_owner_spectator_extension} expresses $V_\xi$ as a core
    with parties in $g(\gamma(C))\not\ni g(p)$ extended by $t^g$, so
    Theorem~\ref{thm:peps_absent_party_spectator_identity} yields~(\ref{eq:peps_all_gate_nonincident}).
    If $C=\varnothing$, $V_\xi$ is a source-free local operation at $*\ne g(p)$
    (Definition~\ref{def:peps_grouped_block_map}) framed by $t^g$, so
    Theorem~\ref{thm:peps_source_free_local_part_restriction} and
    Theorem~\ref{thm:peps_party_restriction} again give~(\ref{eq:peps_all_gate_nonincident}).
\end{proof}

\begin{theorem}[Restriction of a framed source-free composition]
    \label{thm:peps_retained_spectator_framing}
    \lean{TNLean.Lattice.PartyCircuit.Word.restrict_frameList_eval_heq}
    \leanok
    \uses{thm:peps_party_restriction,thm:peps_basis_memory_equalities}
    % Source: 04-compression.tex, lines 233--267, 351--381, 409--434 and 565--588.
    For a source-free chronological word $W:a\to b$, a spectator list $t$, and
    an owner selector $f$ in \cite[proof of Theorem~5.2]{OpenAI2026PolynomialPEPS},
    the framed word $F_t(W)$ satisfies $(F_t(W))|_f=F_{t|_f}(W|_f)$ under the
    retained-list identifications.
\end{theorem}

\begin{proof}\leanok
    \uses{thm:peps_party_restriction,
        thm:peps_operator_composition_framing,thm:peps_basis_memory_equalities}
    By induction on $t$ using Theorem~\ref{thm:peps_party_restriction},
    Theorem~\ref{thm:peps_operator_composition_framing}, and
    Theorem~\ref{thm:peps_basis_memory_equalities}, each retained head register
    frames the inductive restriction while each erased head register is dropped.
\end{proof}

\begin{theorem}[Composition of equal selected local parts]
    \label{thm:peps_selected_local_part_composition}
    \lean{TNLean.Lattice.PartyCircuit.Word.restrict_localPart_comp_eval_heq}
    \leanok
    \uses{thm:peps_source_preparation,thm:peps_party_restriction,
        thm:peps_basis_memory_equalities}
    % Source: 04-compression.tex, lines 246--267, 351--381 and 409--434.
    Let $W:a\to b$, $V:b\to c$, $W':a'\to b'$, $V':b'\to c'$ be chronological
    words with source-register lists $s_W,s_V,s_{W'},s_{V'}$, and let $f$ be an
    owner selector (\cite[proof of Theorem~5.2]{OpenAI2026PolynomialPEPS}). If
    $a|_f=a'|_f$, $b|_f=b'|_f$, $c|_f=c'|_f$, $s_W|_f=s_{W'}|_f$,
    $s_V|_f=s_{V'}|_f$, $(L(W))|_f=(L(W'))|_f$, and $(L(V))|_f=(L(V'))|_f$,
    then $(L(V\circ W))|_f=(L(V'\circ W'))|_f$.
\end{theorem}

\begin{proof}\leanok
    \uses{thm:peps_source_preparation,
        thm:peps_retained_spectator_framing,
        thm:peps_operator_composition_framing,
        thm:peps_identity_input_identification}
    By Theorem~\ref{thm:peps_source_preparation}, $L(V\circ W)=L(V)\circ F_{s_V}(L(W))$
    up to the associative input identification. Restricting via
    Theorem~\ref{thm:peps_retained_spectator_framing},
    Theorem~\ref{thm:peps_operator_composition_framing}, and
    Theorem~\ref{thm:peps_identity_input_identification} gives
    $(L(V\circ W))|_f=(L(V))|_f\circ F_{s_V|_f}((L(W))|_f)=(L(V'\circ W'))|_f$.
\end{proof}

\begin{theorem}[Register moves with equal retained prefixes]
    \label{thm:peps_register_move_retained_prefix}
    \lean{TNLean.Lattice.PartyCircuit.Word.restrict_moveHead_eval_heq_of_restrict_eq}
    \leanok
    \uses{thm:peps_party_restriction,thm:peps_register_exchange,
        thm:peps_basis_memory_equalities}
    % Source: 04-compression.tex, lines 246--267, 351--381 and 409--434.
    For a register $r$, ordered register lists $s,s',a,a'$, and an owner
    selector $f$ with $s|_f=s'|_f$ and $a|_f=a'|_f$
    (\cite[proof of Theorem~5.2]{OpenAI2026PolynomialPEPS}), the register moves
    $M_r(s,a):sra\to rsa$ and $M_r(s',a'):s'ra'\to rs'a'$ satisfy
    $(M_r(s,a))|_f=(M_r(s',a'))|_f$ under the retained-list identifications.
\end{theorem}

\begin{proof}\leanok
    \uses{thm:peps_party_restriction,thm:peps_register_exchange,
        thm:peps_basis_memory_equalities}
    By Theorem~\ref{thm:peps_party_restriction},
    Theorem~\ref{thm:peps_register_exchange}, and
    Theorem~\ref{thm:peps_basis_memory_equalities}, both sides restrict to
    $M_r(s|_f,a|_f)=M_r(s'|_f,a'|_f)$ when $r$ is retained and to
    $\Id_{\mathcal H_{(s|_f)(a|_f)}}=\Id_{\mathcal H_{(s'|_f)(a'|_f)}}$ when $r$ is erased.
\end{proof}

\begin{theorem}[Framing equal selected local parts]
    \label{thm:peps_selected_local_part_framing}
    \lean{TNLean.Lattice.PartyCircuit.Word.restrict_localPart_frame_eval_heq}
    \leanok
    \uses{thm:peps_source_preparation,thm:peps_party_restriction,
        thm:peps_basis_memory_equalities}
    % Source: 04-compression.tex, lines 246--267, 351--381 and 409--434.
    Let $W:a\to b$ and $W':a'\to b'$ be chronological words with source-register
    lists $s_W,s_{W'}$, let $r$ be a register, and let $f$ be an owner selector
    with $a|_f=a'|_f$, $b|_f=b'|_f$, $s_W|_f=s_{W'}|_f$, and
    $(L(W))|_f=(L(W'))|_f$ (\cite[proof of Theorem~5.2]{OpenAI2026PolynomialPEPS}).
    Then
    \begin{align}
        (L(F_r(W)))|_f
         & =(L(F_r(W')))|_f.
        \label{eq:peps_selected_local_part_framing}
    \end{align}
\end{theorem}

\begin{proof}\leanok
    \uses{thm:peps_source_preparation,thm:peps_party_restriction,
        thm:peps_register_move_retained_prefix,
        thm:peps_operator_composition_framing}
    By Theorem~\ref{thm:peps_source_preparation}, $L(F_r(W))=F_r(L(W))\circ M_r(s_W,a)$.
    Here $(M_r(s_W,a))|_f=(M_r(s_{W'},a'))|_f$ by
    Theorem~\ref{thm:peps_register_move_retained_prefix}, and
    $(F_r(L(W)))|_f=(F_r(L(W')))|_f$ by Theorem~\ref{thm:peps_party_restriction}
    and Theorem~\ref{thm:peps_operator_composition_framing}, giving~(\ref{eq:peps_selected_local_part_framing}).
\end{proof}

\begin{theorem}[Local dependence on given incident choices]
    \label{thm:peps_all_branch_incident_locality}
    \lean{TNLean.Lattice.PartyCircuit.SourceCircuit.restrict_partialWord_all_localPart_eval_heq_of_projectAllChoices_eq}
    \leanok
    \uses{def:peps_partial_source_choices,def:peps_touched_labels,
        thm:peps_incident_gate_labels,def:peps_affected_owner,
        def:peps_single_party_selector,thm:peps_source_preparation}
    % Source: 04-compression.tex, lines 246--267, 351--381, 409--434 and 565--588.
    Let $W:a\to b$ be a chronological circuit on $P$, fix $p\in P$, and let
    $g:P\to P\sqcup\{*\}$ be the all-party owner map
    (Definition~\ref{def:peps_affected_owner};
    \cite[proof of Theorem~5.2]{OpenAI2026PolynomialPEPS}). For all-party
    choices $\xi,\zeta\in\Xi_P$ with partial words $V_\xi,V_\zeta$
    (Definition~\ref{def:peps_partial_source_choices}), if
    $\pi_{\{p\}}\xi=\pi_{\{p\}}\zeta$ (Definition~\ref{def:peps_touched_labels}),
    then $(L(V_\xi))|_{g(p)}=(L(V_\zeta))|_{g(p)}$.
\end{theorem}

\begin{proof}\leanok
    \uses{def:peps_partial_source_choices,thm:peps_canonical_source_recovery,
        thm:peps_touched_labels,thm:peps_incident_gate_labels,
        thm:peps_selected_local_part_composition,
        thm:peps_selected_local_part_framing,
        thm:peps_all_gate_nonincident_identity}
    By Theorem~\ref{thm:peps_canonical_source_recovery}, $s_{V_\xi}=s_{V_\zeta}$,
    so $s_{V_\xi}|_{g(p)}=s_{V_\zeta}|_{g(p)}$. We induct on $W$: compositions
    and framed words follow from Theorem~\ref{thm:peps_selected_local_part_composition}
    and Theorem~\ref{thm:peps_selected_local_part_framing}; gates incident to
    $p$ receive the same branch label under $\xi$ and $\zeta$ by
    Theorem~\ref{thm:peps_touched_labels} and Theorem~\ref{thm:peps_incident_gate_labels};
    and gates nonincident to $p$ satisfy
    $(L(V_\xi))|_{g(p)}=\Id=(L(V_\zeta))|_{g(p)}$ by
    Theorem~\ref{thm:peps_all_gate_nonincident_identity}.
\end{proof}

\begin{theorem}[Restriction of a register move]
    \label{thm:peps_register_move_restriction}
    \lean{TNLean.Lattice.PartyCircuit.Word.restrict_moveHead_true_eval_heq,
        TNLean.Lattice.PartyCircuit.Word.restrict_moveHead_false_eval_heq}
    \leanok
    \uses{thm:peps_party_restriction,thm:peps_register_exchange,
        thm:peps_basis_memory_equalities}
    % Source: 04-compression.tex, lines 233--267 and 409--434.
    For a register $r$, ordered register lists $s,t$, and an owner selector $f$
    in \cite[proof of Theorem~5.2]{OpenAI2026PolynomialPEPS}, the register move
    $M_r(s,t):srt\to rst$ satisfies $(M_r(s,t))|_f=M_r(s|_f,t|_f)$ if
    $f(\operatorname{owner}(r))=1$ and
    $(M_r(s,t))|_f=\Id_{\mathcal H_{(s|_f)(t|_f)}}$ if
    $f(\operatorname{owner}(r))=0$.
\end{theorem}

\begin{proof}\leanok
    \uses{thm:peps_party_restriction,thm:peps_register_exchange,
        thm:peps_operator_composition_framing,thm:peps_basis_memory_equalities}
    Induct on $s$ using Theorem~\ref{thm:peps_party_restriction},
    Theorem~\ref{thm:peps_register_exchange},
    Theorem~\ref{thm:peps_operator_composition_framing}, and
    Theorem~\ref{thm:peps_basis_memory_equalities}: for $s=us'$, $M_r(us',t)$
    frames $M_r(s',t)$ by $u$ and then swaps $u$ and $r$, which restricts to
    the retained swap when both $u$ and $r$ are selected and to $\Id$ whenever
    either is erased.
\end{proof}


\begin{theorem}[Erasing a framed register from the local part]
    \label{thm:peps_selected_local_part_erased_frame}
    \lean{TNLean.Lattice.PartyCircuit.Word.restrict_localPart_frame_false_eval_heq}
    \leanok
    \uses{thm:peps_source_preparation,thm:peps_party_restriction}
    % Source: 04-compression.tex, lines 246--267, 351--381 and 409--434.
    For a chronological word $W:a\to b$ with source-register list $s_W$, a
    register $r$, and an owner selector $f$ with $f(\operatorname{owner}(r))=0$
    in \cite[proof of Theorem~5.2]{OpenAI2026PolynomialPEPS}, under the
    identifications $((s_Wra)|_f)=(s_W|_f)(a|_f)$ and $(rb)|_f=b|_f$,
    \begin{align}
        (L(F_r(W)))|_f & =(L(W))|_f.
        \label{eq:peps_selected_local_part_erased_frame}
    \end{align}
\end{theorem}

\begin{proof}\leanok
    \uses{thm:peps_source_preparation,thm:peps_register_move_restriction,
        thm:peps_party_restriction,thm:peps_operator_composition_framing,
        thm:peps_basis_memory_equalities}
    By Theorem~\ref{thm:peps_source_preparation}, $L(F_r(W))=F_r(L(W))\circ M_r(s_W,a)$.
    Since $f(\operatorname{owner}(r))=0$, $(M_r(s_W,a))|_f=\Id$ by
    Theorem~\ref{thm:peps_register_move_restriction} and
    $(F_r(L(W)))|_f=(L(W))|_f$ by Theorem~\ref{thm:peps_party_restriction};
    composing via Theorem~\ref{thm:peps_operator_composition_framing} and
    Theorem~\ref{thm:peps_basis_memory_equalities} gives~(\ref{eq:peps_selected_local_part_erased_frame}).
\end{proof}
